\documentclass[a4paper, serif, 10pt]{moderncv}

%% moderncv
% color and style
\moderncvcolor{black}
\moderncvstyle{banking}

%% page layout/margins
\usepackage[top=2.5cm, bottom=2.5cm, left=1.5cm, right=1.5cm]{geometry}

\nopagenumbers{}

%% Babel config for Italian language
% ref:
% - https://latex3.github.io/babel/guides/locale-italian.html
\usepackage[italian]{babel}

%% fontspec
\usepackage{fontspec}

\IfFontExistsTF{Linux Libertine}{
  \setmainfont[Ligatures={TeX, Common}, Numbers=OldStyle]{Linux Libertine}
}{}
\IfFontExistsTF{Linux Libertine O}{
  \setmainfont[Ligatures={TeX, Common}, Numbers=OldStyle]{Linux Libertine O}
}{}

%% CTeX
% CJK support, used to show CJK characters in the resume
%
% - heading=false: disable ctex heading
% - punct=kaiming: use kaiming punctuations styles for CJK
% - scheme=plain: use plain scheme, do not override \normalsize font size
% - space=auto: space settings for CJK characters
%
% CJK fonts are resolved per locale language, e.g. Noto Serif/Sans CJK JP
% for Japanese resumes, with fallback to the next region if not installed
%
% ref:
% - http://ctan.mirrorcatalogs.com/language/chinese/ctex/ctex.pdf
\usepackage[UTF8, heading=false, punct=kaiming, scheme=plain, space=auto]{ctex}

\IfFontExistsTF{Noto Serif CJK SC}{
  \setCJKmainfont{Noto Serif CJK SC}
}{}
\IfFontExistsTF{Noto Sans CJK SC}{
  \setCJKsansfont{Noto Sans CJK SC}
}{}

%% line spacing
\usepackage{setspace}
\setstretch{1.125}

%% URL styling - use normal text instead of monospace
\urlstyle{same}

%% Patch \httplink and \httpslink to handle full URLs
% The original moderncv macros always prepend http:// or https:// to URLs.
% This causes issues when the URL already has a protocol (e.g., https://example.com)
% resulting in malformed URLs like https://https://example.com.
% This patch checks if the URL already contains "://" and uses it directly if so.
% Note: We use \str_set:Nx (x-expansion) to expand macros like \@homepage first.
\ExplSyntaxOn
\str_new:N \g_yamlresume_url_str

\RenewDocumentCommand{\httpslink}{O{}m}{
  \str_set:Nx \g_yamlresume_url_str {#2}
  \str_if_in:NnTF \g_yamlresume_url_str {://}
    { \url{#2} }
    { \url{https://#2} }
}

\RenewDocumentCommand{\httplink}{O{}m}{
  \str_set:Nx \g_yamlresume_url_str {#2}
  \str_if_in:NnTF \g_yamlresume_url_str {://}
    { \url{#2} }
    { \url{http://#2} }
}
\ExplSyntaxOff

\name{Lorenzo Bianchi}{}
\title{Responsabile Operativo | Logistica e Miglioramento Continuo}
\phone[mobile]{+39 347 123 4567}
\email{lorenzo.bianchi@example.com}
\homepage{https://www.lorenzobianchi.it}

\address{Via Torino 45, Milano, Lombardia, Italia, 20123}{}{}

\extrainfo{{\small \faLinkedin}\ \href{https://www.linkedin.com/in/lorenzo-bianchi-ops}{@lorenzo-bianchi-ops} {} {} {} • {} {} {} 
{\small \faMedium}\ \href{https://medium.com/@lorenzobianchi}{@lorenzobianchi}}

\begin{document}

\maketitle

\section{Informazioni personali}

\cvline{}{\begin{itemize}
\item Responsabile Operativo con oltre 8 anni di esperienza nella gestione di processi logistici, produzione e filiera logistica in contesti industriali e retail.
\item Comprovata capacità di guidare team multidisciplinari, ridurre i costi operativi e migliorare l'efficienza attraverso metodologie Lean Six Sigma e decisioni basate sui dati.
\item Esperienza nella pianificazione della domanda, gestione dei fornitori, ottimizzazione delle scorte e implementazione di sistemi ERP e WMS.
\item Orientato ai risultati, con solide competenze nella gestione del cambiamento, nella negoziazione e nel miglioramento continuo.
\end{itemize}}
\section{Formazione}

\cventry{set 2011–lug 2013}
        {Laurea magistrale, Ingegneria Gestionale, Punteggio: 108/110}
        {Politecnico di Milano}
        {\url{https://www.polimi.it/}}
        {}
        {\begin{itemize}
\item Approfondimento di modelli di pianificazione, controllo di gestione e operations management
\item Progetto di tesi sull'ottimizzazione dei flussi logistici in un'azienda manifatturiera
\item Partecipazione a laboratori su Lean Six Sigma e simulazione di processi produttivi
\end{itemize}
\textbf{Corsi}: Gestione della Produzione e dei Servizi, Logistica e Supply Chain Management, Sistemi Informativi Aziendali, Ricerca Operativa, Economia e Organizzazione Aziendale, Gestione dei Progetti}

\cventry{set 2008–lug 2011}
        {Laurea triennale, Ingegneria Industriale, Punteggio: 105/110}
        {Università di Bologna}
        {\url{https://www.unibo.it/}}
        {}
        {\begin{itemize}
\item Basi solide di analisi quantitativa e gestione industriale
\item Esperienza di tirocinio in un reparto di pianificazione della produzione
\item Sviluppo di competenze nella risoluzione di problemi e nel lavoro di squadra
\end{itemize}
\textbf{Corsi}: Statistica, Ricerca Operativa, Tecnologie Meccaniche, Sistemi di Produzione, Economia Aziendale}
\section{Esperienza professionale}

\cventry{gen 2020–Presente}
        {Responsabile Operativo}
        {Logistica Adriatica S.p.A.}
        {\url{https://www.logisticaadriatica.it}}
        {}
        {\begin{itemize}
\item Direzione di 4 responsabili di area e 45 operatori distribuiti su 3 magazzini, con gestione di un budget operativo annuo di 8 milioni di euro.
\item Riduzione del 18\% dei costi di movimentazione e del 12\% degli scarti tramite progetti Lean Six Sigma e riorganizzazione dei flussi.
\item Implementazione di un nuovo WMS integrato con SAP, con miglioramento della precisione delle scorte dal 92\% al 98,5\%.
\item Definizione di KPI settimanali su produttività, puntualità delle consegne e sicurezza, con reportistica Power BI per il management.
\item Coordinamento di fornitori di trasporto e logistica, con rinegoziazione dei contratti e riduzione del 9\% delle tariffe.
\end{itemize}
\textbf{Parole chiave}: Filiera logistica, Lean Six Sigma, WMS, SAP, KPI, Budget, Sicurezza, Negoziazione}

\cventry{mar 2016–dic 2019}
        {Responsabile di Produzione}
        {Gruppo Meccanica Verona S.r.l.}
        {\url{https://www.gruppomeccanicaverona.it}}
        {}
        {\begin{itemize}
\item Gestione di 2 linee produttive e di un team di 28 persone, con pianificazione della domanda e della capacità produttiva.
\item Introduzione di cicli di pianificazione settimanale e strumenti di visual management, con aumento dell'OEE dal 68\% al 78\%.
\item Riduzione dei fermi macchina del 15\% attraverso manutenzione preventiva e analisi delle cause profonde.
\item Collaborazione con ufficio acquisti e qualità per ridurre i reclami clienti del 20\%.
\end{itemize}
\textbf{Parole chiave}: Produzione, OEE, Manutenzione, Pianificazione, Qualità, Visual Management, Team}

\cventry{giu 2013–feb 2016}
        {Coordinatore Logistico}
        {Retail Italia S.r.l.}
        {\url{https://www.retailitalia.it}}
        {}
        {\begin{itemize}
\item Gestione della distribuzione verso 120 punti vendita, con pianificazione delle rotte e ottimizzazione dei carichi.
\item Riduzione del 14\% dei costi di trasporto tramite consolidamento delle spedizioni e accordi con vettori.
\item Implementazione di procedure di controllo inventariale e riduzione delle differenze di stock dal 3,2\% all'1,1\%.
\end{itemize}
\textbf{Parole chiave}: Distribuzione, Trasporti, Inventario, Ottimizzazione, Retail}
\section{Lingue}

\cvline{Italiano}{Madrelingua o bilingue \hfill \textbf{Parole chiave}: Madrelingua}
\cvline{Inglese}{Competenza professionale completa \hfill \textbf{Parole chiave}: IELTS 8.0, Inglese commerciale}
\cvline{Francese}{Competenza professionale limitata \hfill \textbf{Parole chiave}: DELF B1, Conversazione}
\section{Competenze}

\cvline{Gestione Operativa}{Esperto \hfill \textbf{Parole chiave}: Pianificazione operativa, Gestione budget, KPI, Miglioramento continuo, Gestione fornitori}
\cvline{Filiera Logistica}{Avanzato \hfill \textbf{Parole chiave}: Logistica, Gestione scorte, Pianificazione domanda, Trasporti, Magazzino}
\cvline{Lean Six Sigma}{Avanzato \hfill \textbf{Parole chiave}: DMAIC, Kaizen, 5S, Mappatura flusso valore, Analisi cause profonde}
\cvline{Sistemi e Strumenti}{Avanzato \hfill \textbf{Parole chiave}: SAP, WMS, Power BI, Excel avanzato, Tableau}
\cvline{Leadership}{Avanzato \hfill \textbf{Parole chiave}: Gestione team, Negoziazione, Gestione del cambiamento, Comunicazione}
\section{Premi}

\cventry{nov 2023}
        {Logistica Adriatica S.p.A.}
        {Premio Eccellenza Operativa}
        {}
        {}
        {Riconoscimento per aver guidato il progetto di riorganizzazione del magazzino centrale, con riduzione dei costi del 18\% e miglioramento della puntualità delle consegne al 97\%.}

\cventry{giu 2019}
        {Gruppo Meccanica Verona S.r.l.}
        {Premio Miglior Progetto di Miglioramento}
        {}
        {}
        {Premio per il progetto di manutenzione preventiva che ha ridotto i fermi macchina del 15\% e aumentato l'OEE dal 68\% al 78\%.}
\section{Certificazioni}

\cventry{mar 2018}
        {IASSC}
        {Lean Six Sigma Green Belt}
        {\url{https://www.iassc.org/}}
        {}
        {}

\cventry{giu 2020}
        {ASCM}
        {Certified Supply Chain Professional (CSCP)}
        {\url{https://www.ascm.org/}}
        {}
        {}

\cventry{set 2022}
        {SAP}
        {SAP S/4HANA Supply Chain}
        {\url{https://www.sap.com/training-certification.html}}
        {}
        {}
\section{Pubblicazioni}

\cventry{apr 2022}
        {Digitalizzare la filiera logistica: un modello operativo per le PMI italiane}
        {Logistica Management}
        {\url{https://www.logisticamanagement.it/}}
        {}
        {\begin{itemize}
\item Articolo sull'adozione di WMS e analytics nelle PMI manifatturiere
\item Modello di governance per progetti di digitalizzazione logistica
\item Casi studio di tre aziende del Nord Italia
\end{itemize}}
\section{Progetti}

\cventry{gen 2022–dic 2022}
        {Riorganizzazione dei flussi di approvvigionamento e delle scorte di sicurezza per una rete di tre magazzini.}
        {Progetto Zero Scorte}
        {\url{https://www.logisticaadriatica.it/progetti/zero-scorte}}
        {}
        {\begin{itemize}
\item Analisi ABC/XYZ e ridefinizione dei punti di riordino
\item Riduzione del 22\% del capitale immobilizzato in scorte
\item Miglioramento del livello di servizio dal 94\% al 98\%
\end{itemize}
\textbf{Parole chiave}: Scorte, ABC/XYZ, Livello di servizio, WMS}

\cventry{feb 2021–nov 2021}
        {Ottimizzazione della rete di distribuzione e dei contratti con i vettori per il canale retail.}
        {Rete Distributiva Nord Italia}
        {\url{https://www.logisticaadriatica.it/progetti/rete-nord}}
        {}
        {\begin{itemize}
\item Mappatura dei flussi e ridefinizione delle rotte
\item Riduzione del 12\% dei costi di trasporto
\item Miglioramento della puntualità delle consegne al 96\%
\end{itemize}
\textbf{Parole chiave}: Distribuzione, Trasporti, Retail, Ottimizzazione}
\section{Interessi}

\cvline{Corsa}{Maratona, Trail, Mezza maratona}
\cvline{Vela}{Regate, Costiera, Equipaggio}
\cvline{Lettura}{Saggistica, Economia, Storia}
\section{Volontariato}

\cventry{set 2018–lug 2024}
        {Coordinatore Volontari}
        {Fondazione Banco Alimentare Lombardia}
        {\url{https://www.bancoalimentare.it/}}
        {}
        {\begin{itemize}
\item Coordinamento di 20 volontari durante le raccolte alimentari
\item Pianificazione turni e gestione delle relazioni con i punti vendita
\item Supporto nella logistica di stoccaggio e distribuzione delle derrate
\end{itemize}}
    
\end{document}